\chapter{Modelling Rigid Body Systems}
\label{chap:4}
A rigid-body system is an assembly of component parts. Those components include rigid bodies, joints and various force elements. The joints are responsible for the kinematic constraints in the system, so we interpret the term ‘joint’ broadly to mean any possible kinematic relationship between a pair of rigid bodies. In this chapter, we look at ways of describing a rigid-body system in terms of its component parts. In particular, we seek to describe systems consisting solely of bodies and joints. Such a description can be called a system model, on the grounds that it describes the system itself, rather than some aspect of its behaviour. The role of the system model can be understood as follows. Suppose we wish to calculate the dynamics of a given rigid-body system on a computer. There are basically two ways to proceed: 1. work out the symbolic equations of motion of the given system, and translate these equations into computer code; or 2. construct a system model describing the given system, and supply it as an argument to either (a) a model-based dynamics calculation routine, or (b) a model-based automatic equation (and code) generator. The second approach is preferable, partly because it is much easier to construct and modify a system model than to work out manually the equations of motion, and partly because a model-based routine can be thoroughly debugged and optimized in anticipation of being used a large number of times. The dynamics algorithms in this book are model-based, and expect to be given a system model along the lines described in this chapter. Approach 2(b) is considered in Section 10.4.

base J4 B3 J3 B2 J2 B1 J1 base B2 B1 B3 J1 J2 J3 J4 6 DoF body dish body dish planet 6 DoF Figure 4.1: Connectivity graphs of two rigid-body systems 4.1 Connectivity A system model must contain the following data: a description of the component parts, and a description of how they are connected together. The latter is referred to as the system’s topology, or connectivity, and it can be expressed in the form of a connectivity graph having the following properties: 1. The nodes represent bodies. 2. The arcs represent joints. 3. Exactly one node represents a fixed body, which is known as the base. All other nodes represent moving bodies. 4. The graph is undirected. 5. The graph is connected (i.e., a path exists between any two nodes). The term ‘moving body’ includes any body that is capable of moving, whether or not it actually moves. Thus, a body that has been immobilized by kinematic constraints would count as a moving body. From here on, the term and ‘graph’ should be taken to mean ‘undirected graph’, unless indicated otherwise. Some examples of connectivity graphs are shown in Figure 4.1. The first is a simple planar mechanism called a crank-slider linkage. It consists of three moving bodies (B1 to B3), a fixed base, three revolute joints (J1 to J3) and one prismatic joint (J4). Rotation of body B1 (the crank) causes B3 to slide back and forth relative to the base. In the connectivity graph, each node represents

one body, and each arc one joint. It is usually necessary to label at least some of the nodes and arcs, so that the correspondence between the graph and the original mechanism is clear. In addition, it is sometimes useful to make the base node stand out from the others, for easy visual identification. We accomplish this by marking the base node with a white dot. Note that the bodies in a mechanism are usually called links. The second example is a satellite in orbit. The satellite itself consists of two moving bodies and one joint. However, we cannot construct a connectivity graph directly from this two-body system, as rule 3 requires that one node must represent a fixed base. We therefore add a third body—in this case, the planet below—to serve as a fixed base. Having added a third body, rule 5 requires that we also add a second joint, in order to connect the planet to one of the other two bodies. As there is no physical connection between the satellite and the planet, we use a special joint that allows six degrees of motion freedom—a 6-DoF joint. This joint imposes no kinematic constraints, and therefore has no physical effect on the system. Its role is to define additional joint variables, so that the position and velocity variables for the system as a whole can be drawn from the position and velocity variables of the joints it contains. In this case, the satellite has a total of 6 + ns degrees of freedom, where ns is the degree of freedom allowed by the joint on the satellite. We therefore need a total of 6+ns velocity variables, and that is exactly the number defined by the two joints. A body that is connected to a fixed base via a 6-DoF joint is called a floating base. In this example, the satellite’s body was chosen as the floating base, but we could have chosen the dish instead. 4.1.1 Kinematic Trees and Loops A graph is a topological tree if there exists exactly one path between any two nodes in the graph. If the connectivity graph of a rigid-body system is a topological tree, then we call the system itself a kinematic tree. The satellite in Figure 4.1 is a kinematic tree, but the crank-slider linkage is not. Kinematic trees have special properties that make it easy to calculate their dynamics efficiently. Let G be any connectivity graph. A spanning tree of G, denoted Gt, is a subgraph of G containing all of the nodes in G, together with any subset of the arcs in G such that Gt is a topological tree. Every connectivity graph has at least one spanning tree. If G itself is a tree, then Gt = G; otherwise, Gt is a proper subgraph of G and is not unique. Figure 4.2 shows some examples of spanning trees. A topological tree with n nodes has n −1 arcs. Therefore, if G contains nn nodes and na arcs, then Gt will contain nn nodes and nn −1 arcs, implying that there will be na −nn + 1 arcs that are in G but not in Gt. These non-tree arcs are called chords. Although the number of chords in G is fixed, the set of arcs that happen to be chords will vary with the choice of Gt. This variation is illustrated in Figure 4.2. (We distinguish chords from tree arcs by drawing them with dashed lines.)

every possible spanning tree original graph chords and their cycles chord cycle Figure 4.2: Spanning trees, chords and cycles A cycle is a closed path in a graph. There are no cycles in a tree, so every cycle in a graph must traverse at least one chord. We are particularly interested in those cycles that traverse exactly one chord. There is always exactly one such cycle for each chord in the graph, but its route depends on the choice of spanning tree. For example, the graph in Figure 4.2 has two chords and three cycles. For any choice of spanning tree, two of these cycles are one-chord cycles, and the third (which is not shown) traverses both chords. The path of a one-chord cycle consists of the one chord it traverses, together with the unique path in the spanning tree between the two nodes connected by the chord. Each cycle in a connectivity graph defines a kinematic loop. However, we are not interested in them all, but only in the independent kinematic loops. This is the minimal set of loops that must be taken into account when formulating the equations of motion. The importance of the one-chord cycles is that they identify the set of independent kinematic loops. If we describe a rigid-body system as having n kinematic loops, then what we really mean is that it has n independent kinematic loops. Consider a rigid-body system comprising a fixed base, NB moving bodies and NJ joints, the latter including any 6-DoF joints that were added to ensure that the connectivity graph was connected. Let G be the connectivity graph of this system, and let Gt be a spanning tree of G. Both graphs will contain NB +1 nodes. However, G will contain NJ arcs, whereas Gt will contain only NB, so the number of chords in the system is NJ −NB. The number of independent kinematic loops in the system, NL, equals the number of chords in G, and is therefore given by the formula NL = NJ −NB . (4.1)

Applying this formula to the crank-slider linkage in Figure 4.1 reveals that it has one kinematic loop. If a system has kinematic loops, then it is called a closed-loop system, otherwise it is a kinematic tree. The joints associated with the chords are known as loop joints, or loop-closing joints, and those that are part of the spanning tree are tree joints. A rigid-body system therefore has NB tree joints and NL loop joints. 4.1.2 Numbering Bodies and joints can be identified in various ways, but the most convenient way is to number them. We shall use a numbering scheme that is well suited to describing and implementing dynamics algorithms. According to this scheme, the bodies are numbered from 0 to NB, the joints from 1 to NJ, and the numbers are assigned according to the following rules. 1. Choose a spanning tree, Gt. 2. Assign the number 0 to the node representing the fixed base, and define this node to be the root node of Gt. 3. Number the remaining nodes from 1 to NB in any order such that each node has a higher number than its parent in Gt. 4. Number the arcs in Gt from 1 to NB such that arc i connects between node i and its parent. 5. Number all remaining arcs from NB + 1 to NJ in any order. 6. Each body gets the same number as its node, and each joint gets the same number as its arc. This scheme is called regular numbering. If these rules are applied to an unbranched kinematic tree (a tree in which no node has more than one child), then the bodies and joints are numbered consecutively from the base, as shown in Figure 4.3. In all other cases, the numbering is not unique. For example, Figure 4.4 shows every possible numbering of the crank-slider linkage in Figure 4.1. These eight possibilities arise from four possible spanning trees. Two trees are unbranched, and therefore have only one numbering. These two appear in the top left and bottom right graphs in the figure. The other two are branched at the root node, and they have three numberings each. 4.1.3 Joint Polarity Each joint connects between two bodies. We identify one of these bodies as the joint’s predecessor, and the other as the joint’s successor. The joint itself is said to connect from its predecessor to its successor. The purpose of this terminology

1 2 3 0 1 3 2 n−2 n−1 n−1 n n Figure 4.3: Unique numbering of an unbranched kinematic tree 1 2 3 4 0 1 3 2 1 2 3 4 0 1 3 2 1 2 3 4 0 1 3 2 1 2 3 4 0 1 3 2 1 2 3 4 0 1 3 2 1 2 3 4 0 1 3 2 1 2 3 4 0 1 3 2 0 2 3 3 4 1 1 2 Figure 4.4: Every possible numbering of the crank-slider linkage in Figure 4.1 is twofold: to define the relationship between joint and body variables, and to specify which way round the joint is connected in the system. For example, the velocity across a joint is defined to be the velocity of the successor relative to the predecessor, and the force transmitted across a joint is defined to be the force transmitted from the predecessor to the successor. The identification of which way round a joint is connected between two bodies is called its polarity. We can identify the polarity of a joint in a connectivity graph by placing an arrow on the arc such that it points from the predecessor to the successor. Examples of this are shown in Figures 4.5 and 4.6. As a matter of convenience, we adopt the following convention on the polarity of tree joints: the normal polarity for a tree joint is to connect from the parent to the child. All tree joints are assumed to have this polarity unless indicated otherwise. Some joints are physically symmetrical, meaning that the motion of the successor relative to the predecessor is qualitatively the same as the motion of the predecessor relative to the successor. Joints that do not have this property are asymmetrical. Switching the polarity of a symmetrical joint does not make a material difference to a rigid-body system, but switching the polarity of an asymmetrical joint does. For example, if two bodies are connected by a balland-socket joint, which is symmetrical, then it does not really matter which body has the ball and which the socket—if we were to switch them round, then the motion would still be the same. In contrast, if the two bodies are connected by a rack-and-pinion joint, which is asymmetrical, then it does matter which body has the rack and which the pinion. Thus, the polarity of an asymmetrical

1 2 3 4 0 1 3 2 p = [0, 1, 2, 3] s = [1, 2, 3, 0] 1 2 3 4 0 1 3 2 p = [0, 0, 1, 2] s = [1, 2, 3, 3] Figure 4.5: Joint predecessor and successor arrays 0 1 2 3 4 5 6 7 8 9 10 11 p = [0, 1, 1, 4, 3, 4, 4, 4, 2, 6, 5] s = [1, 2, 3, 0, 5, 6, 7, 8, 5, 3, 7] λ = [0, 1, 1, 0, 3, 4, 4, 4] Figure 4.6: Predecessor, successor and parent arrays joint is determined by the physical system, but the polarity of a symmetrical joint is not, leaving us free to choose its polarity in the system model. 4.1.4 Representation There are many ways to represent a connectivity graph in a computer. One of the simplest is to record the joint predecessor and successor body numbers in a pair of arrays, p and s, such that p(i) and s(i) are the predecessor and successor body numbers, respectively, of joint i. This method is illustrated in Figure 4.5, which reproduces two of the connectivity graphs from Figure 4.4, adding arrows to show the polarity of each joint. Together, p and s provide a complete description of a connectivity graph, from which other useful quantities can be calculated. The first of these is the parent array, λ, which identifies the parent of each body in the spanning tree: λ(i) is the parent of body i. According to rules 3 and 4 of the numbering scheme, tree joint i connects between bodies i and λ(i), and λ(i) < i, so λ(i) = min(p(i), s(i)) . (1 ≤i ≤NB) (4.2) The parent arrays for the two graphs in Figure 4.5 are λ = [0, 1, 2] and λ = [0, 0, 1]; and a larger example is shown in Figure 4.6. If a tree joint satisfies s(i) = i, then it connects from the parent to the child, and is said to be in the forward direction in the tree. Otherwise, it connects from the child to the parent, and is said to be in the reverse direction. The tree joints in Figure 4.5 are all in the forward direction. (The arrows point away from the root.) The tree joints in Figure 4.6 are also all in the forward direction, except for joint 4. The tree-joint numbers and forward-pointing arrows have been

omitted from this diagram: the former because they can be deduced from the numbering scheme, and the latter because we regard ‘forward’ as the standard direction for a tree joint. Three more quantities provide useful information about the connectivity of a kinematic tree: for any body i except the base, κ(i) is the set of all tree joints on the path between body i and the base; and for any body i, including the base, µ(i) is the set of children of body i, and ν(i) is the set of bodies in the subtree starting at body i. They are called the support, child and subtree sets, respectively. (A tree joint is said to support a body if it lies on the path between that body and the base.) ν(i) can also be regarded as the set of all bodies supported by joint i. For the spanning tree in Figure 4.6, the support, child and subtree sets are as follows: κ(1) = \{1\} µ(0) = \{1, 4\} ν(0) = \{0, 1, 2, . . ., 8\} κ(2) = \{1, 2\} µ(1) = \{2, 3\} ν(1) = \{1, 2, 3, 5\} κ(3) = \{1, 3\} µ(2) = \{\} ν(2) = \{2\} ... ... ... κ(8) = \{4, 8\} µ(8) = \{\} ν(8) = \{8\} κ, λ, µ and ν obviously have a lot of simple properties that follow from their definitions. For example, µ(i) ⊆ν(i), j∈κ(i) ⇒i∈ν(j), and so on. One not-so-obvious property that turns out to be quite useful is the identity NB X i=1 X j∈κ(i) ( · · · ) = NB X j=1 X i∈ν(j) ( · · · ) . (4.3) The left-hand side is a summation over all i, j pairs in which joint j supports body i, but the right-hand side is also a summation over all i, j pairs in which joint j supports body i, and so the two are the same. Example 4.1 The role of λ in dynamics algorithms is illustrated by the following simple example. Suppose we wish to calculate the velocity of every body in a rigid-body system, given the velocities across the tree joints. Let vi and vJi be the velocity of body i and the velocity across joint i, respectively, the latter being the velocity of the child relative to the parent (i.e., the joints are in the forward direction). Now, the velocity of body i can be defined recursively as the sum of its parent’s velocity and the velocity across the joint connecting it to its parent; so the formula for calculating the body velocities is vi = vλ(i) + vJi . (v0 = 0) The algorithm for calculating these velocities is therefore v0 = 0 for i = 1 to NB do vi = vλ(i) + vJi end

The property λ(i) < i ensures that vλ(i) is always calculated before it is used. Another way to express the body velocities is vi = X j∈κ(i) vJj . (4.4) Example 4.2 Example 2.3 introduced a matrix called the body Jacobian, which gives the spatial velocity of an individual body as a function of the joint velocity variables. If Ji is the body Jacobian for body i, then vi = Ji ˙q . Equation 2.18 gave a formula for Ji that is valid for the special case of an unbranched kinematic chain, so let us now obtain the general formula. Let the velocity across joint j be Sj ˙qj, and let ϵij be defined as follows: ϵij =  1 if j ∈κ(i) 0 otherwise. (4.5) So ϵij is nonzero if joint j supports body i. The velocity of body i can now be expressed as follows: vi = X j∈κ(i) Sj ˙qj = NB X j=1 ϵijSj ˙qj = ϵi1S1 · · · ϵiNBSNB 

˙q1 ... ˙qNB

. Thus, the general formula for the body Jacobian is Ji = ϵi1S1 ϵi2S2 · · · ϵiNBSNB  . (4.6) 4.2 Geometry If two bodies are connected by a joint, then their relative motion is constrained. A complete description of the constraint requires knowledge of two things: the motion allowed by the joint, and the location of the joint relative to each body. The former is described by a joint model. The latter is described by the system’s geometry. A geometric model of a rigid-body system is a description of where the connections are. In particular, it is a description of the relative position of each joint on each rigid body. Figure 4.7 shows a geometric model for a rigid-body system comprising a fixed base, B0, three moving bodies, B1 to B3, three tree joints, J1 to J3, and one loop joint, J4. The dashed lines in this figure indicate only the connectivity of each joint, not its exact physical location. Each moving body has a coordinate frame embedded in it, which defines a local coordinate system for that body. These frames are labelled F1 to F3,

F0 F1 F2 F3 B1 B2 B3 F0,1 F1,2 F1,3 F2,4 F3,4 XJ1 XJ2 XJ3 B0 = base XJ4 0,1X0 1,2X1 1,3X1 2,4X2 3,4X3 J2 J4 J3 J1 Figure 4.7: Geometric model of a rigid-body system and the coordinate systems they define are known as body coordinates, or link coordinates. There is also a frame F0, embedded in the base, which defines an absolute, reference, or base coordinate system for the whole rigid-body system. The position of body Bi is defined as the position of Fi relative to F0, and it can be represented by the coordinate transform iX0. The figure also shows several frames with names of the form Fi,j where i refers to a body and j to a joint. These frames identify where the joints are located on each body. Each tree joint i connects from frame Fλ(i),i to frame Fi (or vice versa if the joint is in the reverse direction), and each loop joint i connects from Fp(i),i to Fs(i),i. There are therefore NB + 2NL locating frames in total: one for each tree joint and two for each loop joint. The position of frame Fi,j relative to Fi is given by the transform i,jXi, which is a constant. These transforms, or a set of data from which they can be calculated, define the geometric model of a rigid-body system. For convenience, we define an array of tree transforms, XT, such that XT(i) = λ(i),iXλ(i). These are the locating transforms for the tree joints, and this array provides a complete description of the geometry of the spanning tree. We also define the arrays XP and XS, each indexed from 1 to NL, such that XP(i−NB) = p(i),iXp(i) and XS(i−NB) = s(i),iXs(i) for each loop joint i. These are the locating transforms for the loop joints. Finally, each joint defines a joint transform, XJi, which is a function of the joint’s position variables. These transforms are discussed in Section 4.4. As shown in the figure, XJi is equal to iXλ(i),i if i is a tree joint, and s(i),iXp(i),i if i is a loop joint.

Example 4.3 Suppose we wish to calculate the position of each body in a kinematic tree; that is, we wish to calculate iX0 for each body. The algorithm to accomplish this is for i = 1 to NB do XJ = jcalc(jtype(i), qi) iXλ(i) = XJ XT(i) if λ(i)̸ = 0 then iX0 = iXλ(i) λ(i)X0 end end In this code fragment, XJ is a local variable, jtype(i) is a data structure describing joint i, and jcalc is a function that calculates a joint transform from a joint description and a vector of joint position variables. jtype(i) and jcalc are covered in Section 4.4. 4.3 Denavit-Hartenberg Parameters In general, six parameters are needed to specify the position of one coordinate frame relative to another. However, for certain kinds of rigid-body system, it is possible to locate the body coordinate frames in such a way that fewer parameters are needed. One such method, originally described by Denavit and Hartenberg (Denavit and Hartenberg, 1955), needs only four parameters per joint. It exploits the fact that some of the most common joint types can be characterized by a line in space, and that only four parameters are needed to locate a line. Figure 4.8 shows how the Denavit-Hartenberg (DH) method is applied to an unbranched kinematic tree representing a robot arm or similar device. The important geometrical features of this system are: • a base coordinate frame, • n joint axes, and • an end-effector coordinate frame, which is embedded in the final link. The end-effector frame is needed for the correct execution of positioning commands by the robot, but is not needed for dynamics calculations. The DH coordinate frames are placed according to the following rules. 1. Axes z0 and zn+1 are aligned with the z axes of the base and end-effector coordinate frames, respectively. 2. Axes z1 to zn are aligned with the n joint axes such that zi is aligned with the axis of joint i. (If joint i happens to be prismatic, then zi can be any axis with the correct direction; but it is sensible to choose an axis that intersects with zi−1.)

prismatic joint motion revolute joint motion z0 = zb z1 zn zn+1 x0 x1 xn d0 d1 dn dn+1 θ0 θ1 θn θn+1 a1 a2 an+1 α1 α2 αn+1 end effector coordinate frame base coordinate frame joint 1 axis joint 2 axis joint n axis xb yb xe ye = ze Figure 4.8: Denavit-Hartenberg Coordinate Frames and Parameters

3. Axis xi is the common perpendicular between zi and zi+1, directed from zi to zi+1. If zi and zi+1 intersect, then there is a 180◦ambiguity in the direction of xi. If zi and zi+1 are parallel, then the location of the common perpendicular is indeterminate, and it is reasonable to choose one that intersects with xi−1. 4. The body coordinate frame Fi is defined by the axes xi, yi and zi, where yi is derived from xi and zi to make a right-handed coordinate frame. Having placed the coordinate frames, their relative locations are described by the following DH parameters: di the distance from xi−1 to xi measured along zi. θi the angle from xi−1 to xi measured about zi. ai the distance from zi−1 to zi measured along xi−1. The definition of xi−1 implies ai ≥0. αi the angle from zi−1 to zi measured about xi−1. All of the parameters ai and αi are constants, but some of the parameters di and θi serve as joint variables. Specifically, if joint i is revolute, then θi is its joint variable; if joint i is prismatic, then di is its joint variable; if joint i is cylindrical, then both θi and di are joint variables; and if joint i is a screw joint with a pitch of hi, then θi is the joint variable, and the distance from xi−1 to xi is defined to be di + hi θi. If every joint is revolute, then the DH parameters define the following geometric model for the unbranched tree in Figure 4.8: i = 1 : XT(1) = xlt([a1 0 d1]T) rotx(α1) rotz(θ0) xlt([0 0 d0]T) i > 1 : XT(i) = xlt([ai 0 di]T) rotx(αi) (See Table 2.2 for definitions of of the translation and rotation functions. Also, see §A.4 in the appendix for screw transforms.) It requires a total of 4n+6 DH parameters to describe a system with n joints, of which at least n will be joint variables rather than parameters. However, the last four parameters serve only to locate the end-effector frame, and can be omitted if this information is not needed. Up to five more parameters can be omitted if the base frame can be chosen to be aligned with z1. In its basic form, the Denavit-Hartenberg method is applicable to any rigidbody system in which every joint can be characterized by a joint axis and the topology is either an unbranched kinematic tree or a single closed kinematic loop. However, the method can be adapted to a broader class of systems, as described in Khalil and Dombre (2002). DH parameters are described in various textbooks, such as Angeles (2003) and Craig (2005). The notational details tend to vary from one author to the next, especially the numbering of xi, ai and αi. The notation used here follows the Springer Handbook of Robotics (Siciliano and Khatib, 2008).

q1 Fp Fs x y q2 q3 planar joint q2 q1 Fp Fs x z cylindrical joint Fp Fs 6DoF joint r = (q5,q6,q7) (q1,q2,q3,q4) y (in Fs coords) Figure 4.9: Geometrical interpretation of some joint position variables 4.4 Joint Models A joint can be regarded as a motion constraint between two Cartesian frames— one embedded in the joint’s successor body and the other in its predecessor. We shall call these frames Fs and Fp, respectively. The purpose of a joint model is to provide the information required to calculate the motion of Fs relative to Fp as a function of the joint variables. The kind of motion allowed by a joint depends on its joint type. For example, a revolute joint allows pure rotation about a single axis, whereas a spherical joint allows arbitrary rotation about a specific point. We therefore need one joint model for each type of joint. The contents of a joint model are a collection of data and formulae for calculating specific quantities, such as: the joint transform (sXp), the joint velocity (vJ) the joint’s motion subspace (S), and so on. Formulae for several common joint types are shown in Table 4.1, and the geometries of some of these joints are shown in Figure 4.9. Let us examine one model in more detail. The first entry in Table 4.1 describes a revolute joint that allows Fs to rotate relative to Fp about their common z axis, and defines the joint variable to be the angle of rotation. The two columns marked E and r describe the rotational and translational components of the joint transform, and the transform itself is given by the equation sXp = rot(E) xlt(r). (See Table 2.2 for the definitions of rot, xlt and rz.) Observe that Fs coincides with Fp when the joint variable (q1) is zero. We apply this convention as far as possible to all joint models. Another convention we apply to all joint models is that S, T , and related quantities, are expressed in Fs coordinates. In the case of a revolute joint this does not really matter, since sS = pS, but it does matter for other joints. In a well-designed dynamics simulator, there would likely be a joint model library, containing one model for each type of joint supported by the software. A system model would then contain a joint type descriptor for each joint in the system. We shall use the expression jtype(i) to denote the joint type descriptor of joint i. In the simplest case, a joint type descriptor is just a type code identifying a particular model in the library. For example, the type code ‘R’

Joint Type Joint Transform (sXp) Motion Constraint Force E r Subsp. (S) Subspace (T ) Revolute (hinge joint) rz(q1)

0 0 0

2 6666664 0 0 1 0 0 0 3 7777775 2 6666664 1 0 0 0 0 0 1 0 0 0 0 0 0 0 0 0 0 1 0 0 0 0 0 1 0 0 0 0 0 1 3 7777775 Prismatic (sliding joint) 13×3

0 0 q1

2 6666664 0 0 0 0 0 1 3 7777775 2 6666664 1 0 0 0 0 0 1 0 0 0 0 0 1 0 0 0 0 0 1 0 0 0 0 0 1 0 0 0 0 0 3 7777775 Helical (screw joint) with pitch h rz(q1)

0 0 h q1

2 6666664 0 0 1 0 0 h 3 7777775 2 6666664 1 0 0 0 0 0 1 0 0 0 0 0 0 0 −h 0 0 1 0 0 0 0 0 1 0 0 0 0 0 1 3 7777775 Cylindrical rz(q1)

0 0 q2

2 6666664 0 0 0 0 1 0 0 0 0 0 0 1 3 7777775 2 6666664 1 0 0 0 0 1 0 0 0 0 0 0 0 0 1 0 0 0 0 1 0 0 0 0 3 7777775 Planar (see Fig. 4.9) rz(q1)

c1q2 −s1q3 s1q2 + c1q3 0

2 6666664 0 0 0 0 0 0 1 0 0 0 1 0 0 0 1 0 0 0 3 7777775 2 6666664 1 0 0 0 1 0 0 0 0 0 0 0 0 0 0 0 0 1 3 7777775 Spherical (ball \& socket)  see Eq. 4.12 

0 0 0

2 6666664 1 0 0 0 1 0 0 0 1 0 0 0 0 0 0 0 0 0 3 7777775 2 6666664 0 0 0 0 0 0 0 0 0 1 0 0 0 1 0 0 0 1 3 7777775 6-DoF Joint (free motion)  see Eq. 4.12  E−1

q5 q6 q7

16×6 Notes: c1 = cos(q1), s1 = sin(q1), and sXp = rot(E) xlt(r) (see Table 2.2). Table 4.1: Joint Model Formulae

could identify a revolute joint and ‘P’ a prismatic one. However, some joints are characterized by both a type and a set of parameters. For example, a helical joint is characterized both by the fact that it is helical and a parameter that specifies the pitch of the helix. Each individual helical joint in a rigid-body system can have a different pitch; so parameters like this are a property of individual joints, rather than the joint model. To accommodate joint parameters, we define a joint type descriptor to be a data structure containing both a type code and the set of parameter values (if any) required for a joint of that type. Joint Model Calculation Routine Subsequent chapters describe several dynamics algorithms, all of which make use of joint model data. To keep the algorithm descriptions simple, we assume the existence of a single function, called jcalc, that calculates all requested joint-related data in a single call. A specimen call to this function looks like this: [XJ, S, vJ, cJ] = jcalc(type, q, α) (explicit joint constraint) [δ, T ] = jcalc(type, sXp) (implicit joint constraint) [ ˙q] = jcalc(type, q, α) (for qdfn (Eq. 3.7)) In each case, the expression on the left is a list of requested return values, and the argument type refers to a joint type descriptor (e.g. jtype(i)). The use of jcalc to calculate implicit joint constraints is covered in Chapter 8, and its use to calculate ˙q is covered later in this chapter. In the algorithm listings, we generally assume that α = ˙q, and use ˙q in place of α in the argument lists. If a joint depends explicitly on time, then time must be added as a final argument in the list. Note that vJ, cJ, S and T are all expressed in Fs coordinates. Expressions for vJ and cJ are given in Section 3.5. Example 4.4 Example 4.3 showed how to calculate the coordinate transform iX0 for each body in a kinematic tree. We now extend this example to calculate both iX0 and 0vi (the velocity of body i expressed in base coordinates). The algorithm to accomplish this is 0v0 = 0 for i = 1 to NB do [XJ, vJ] = jcalc(jtype(i), qi, ˙qi) iXλ(i) = XJ XT(i) if λ(i)̸ = 0 then iX0 = iXλ(i) λ(i)X0 end 0vi = 0vλ(i) + 0XivJ end

This code fragment contains a new feature: it calculates iX0, but then uses 0Xi to transform vJ. The explanation for this practice can be found in Appendix A, where it is shown how to calculate X−1v directly from X and v without having to calculate X−1 first. It is also possible to calculate X∗f and (X∗)−1f directly from X and f without having to calculate X∗or (X∗)−1 first. We exploit this fact in several dynamics algorithm listings. Example 4.5 A peculiar feature of both the planar joint and the 6-DoF joint in Table 4.1 is that the translational joint variables express the translation in Fs coordinates rather than Fp coordinates. The reason for doing this is to make the motion subspace matrix a simple constant, so that ˚ S = 0, cJ = 0, and the simple form of S allows various optimizations to be applied to the dynamics algorithms. However, there is potentially a big disadvantage in this arrangement: if the successor body is far distant from the predecessor, and starts to spin, then the spinning will cause large rates of change in the position variables. (Just think of a tumbling satellite or aircraft.) One way to get the best of both worlds is to use a set of position variables in which the translations are expressed in Fp coordinates, together with a set of velocity variables in which the linear velocity is expressed in Fs coordinates. Let pq1, pq2 and pq3 be a set of position variables for a planar joint, such that pq2 and pq3 are translations in the x and y directions of Fp; and let sq1, sq2 and sq3 be a second set of position variables for the same planar joint, such that sq2 and sq3 are translations in the x and y directions of Fs. This second set of variables is the set used in Table 4.1 and shown in Figure 4.9. The relationship between these two sets is

pq1 pq2 pq3

=

1 0 0 0 c1 −s1 0 s1 c1

sq1 sq2 sq3

. We wish to use pqi as the position variables, but s ˙qi as the velocity variables. To cater for this, we need a formula to calculate the derivatives of the position variables from the position and velocity variables. Such a formula is obtained by differentiating the above equation, to get

˙q1 ˙q2 ˙q3

=

1 0 0 −q3 c1 −s1 q2 s1 c1

α1 α2 α3

, where we have dropped the superscript p from pqi and p ˙qi, and substituted αi for s ˙qi. This is the calculation that jcalc must perform if it is called upon to calculate ˙q from q and α for this particular model of the planar joint. Example 4.6 Let us construct a joint model for the rack-and-pinion joint shown in Figure 4.10. This joint allows a pinion that is part of the successor body to roll along a rack that is part of the predecessor body. The rack

r q1 Fp Fs rq1 x y Figure 4.10: Rack-and-pinion joint is parallel to the x axis of Fp, and the pinion is centred on the z axis of Fs. The pinion has a radius of r, which is a parameter of the joint model. The joint allows only a single degree of motion freedom, so there is only one joint variable, denoted q1, and it is chosen to be the rotation angle of the pinion. If Fp and Fs coincide when q1 = 0, then the transform from Fp to Fs consists of a translation by rq1 in the x direction followed by a rotation of q1 about the z axis. The joint transform is therefore XJ = sXp = E 0 0 E   1 0 −p× 1  , where E = rz(q1) =

c1 s1 0 −s1 c1 0 0 0 1

and p =

rq1 0 0

. The next step is to find formulae for vJ and S. We begin by expressing the joint velocity as a function of the joint’s position and velocity variables. Relative to Fp, the pinion is rotating with an angular velocity of ω = [0 0 ˙q1]T about an axis passing through p, and it is translating with a linear velocity of ˙p = [r ˙q1 0 0]T. This equates to a spatial velocity, expressed in Fp coordinates, of pvJ =  ω p × ω  + 0 ˙p  =

0 0 ˙q1

rq1 0 0

×

0 0 ˙q1

+

0 0 0 r ˙q1 0 0

=

0 0 ˙q1 r ˙q1 −rq1 ˙q1 0

. As there is only a single joint variable, the motion subspace matrix is given by

the formula S = vJ/ ˙q1. Therefore vJ = sXp pv =

0 0 ˙q1 rc1 ˙q1 −rs1 ˙q1 0

and S = vJ ˙q1 =

0 0 1 rc1 −rs1 0

. Finally, as S is not a constant, we need a formula for the velocity-product term cJ. This is given by cJ = ˚ S ˙q1 = ∂S ∂q1 ˙q2 1 =

0 0 0 −rs1 ˙q2 1 −rc1 ˙q2 1 0

. (See Eq. 3.42 for ˚ S.) Observe that the rack-and-pinion joint simplifies to a revolute joint if r = 0. Polarity Reversal As explained in Section 4.1.4, a joint in a kinematic tree is said to be in the forward direction if the predecessor is the parent of the successor, and in the reverse direction otherwise. If a joint is symmetrical, then its polarity can be chosen at the time the system model is being defined so that it is in the forward direction; but if it is asymmetrical, then its polarity is determined by the rigidbody system being modelled. Thus, it is not possible to guarantee that every tree joint is in the forward direction. The simplest way to deal with a reverse-direction joint is to convert it to a forward-direction joint. We can do this as follows. First, add a boolean variable to the joint descriptor data structure that takes the value true if the joint is a tree joint in the reverse direction, and takes the value false otherwise. Given this extra piece of information, jcalc can be modified to simulate a joint of the opposite polarity whenever this boolean variable is true. In effect, jcalc reverses the roles of the predecessor and successor. To accomplish this task, jcalc merely has to modify some of the calculated values before they are returned. For example, instead of returning the matrix sXp that has been calculated according to the formula for the given joint type, it returns the inverse of this matrix; and, instead of returning the calculated value of vJ, it returns −pXsvJ. A complete set of modified return values is shown in Table 4.2. The advantage of this approach is that it allows the dynamics algorithms to assume that all tree joints are in the forward direction. This, in turn, makes the algorithms simpler and easier to implement.

normal simulating return value opposite polarity sXp pXs S −pXsS T −pXsT vJ −pXsvJ cJ −pXscJ ˙q ˙q Table 4.2: How jcalc simulates an opposite-polarity joint 4.5 Spherical Motion There are many ways to represent a general 3D rotation (e.g. see Rooney, 1977), but the two main choices are Euler angles and Euler parameters. The latter are also known as unit quaternions. Euler angles have the advantage that the position variables are a set of three independent angles, and the velocity variables are their derivatives. This makes them relatively easy to implement. Their disadvantage is that they suffer from singularities: as the orientation approaches a singularity, the representation becomes increasingly ill-conditioned, and eventually breaks down altogether. The main advantage of Euler parameters is that they are free of singularities; but their disadvantage is that they are a set of four non-independent position variables, which implies that the velocity variables cannot be the derivatives of the position variables. Software that uses Euler parameters must be written to allow the number of position variables to be different from the number of velocity variables. Euler Angles We use the term ‘Euler angles’ in both a broad sense and a narrow sense. In the broad sense, it refers to any method for representing a general rotation by means of a sequence of three rotations about specified coordinate axes. In the narrow sense, it refers to the special case in which the first and third rotations take place about the z axis, and the second rotation takes place about either the x axis or the y axis. Every version of Euler angles suffers from singularities. They occur whenever the first and third axes are aligned. If the first and third rotations use the same coordinate axis, then the singularities occur when the second angle is either 0◦ or 180◦; otherwise, they occur when the second angle is ±90◦. In the vicinity of a singularity, small changes in the rotation being represented can induce large changes in the first and third angles. A joint model based on Euler angles is not valid at a singularity. Let us construct a spherical joint model using yaw, pitch and roll angles. These are a sequence of rotations about the z, y and x axes, in that order.

Thus, the correct orientation of Fs relative to Fp is obtained by first placing Fs coincident with Fp, then rotating it about its z axis by an amount q1 (yaw), then rotating it about its y axis by an amount q2 (pitch), and then rotating it about its x axis by an amount q3 (roll). The formula for E is therefore E = rx(q3) ry(q2) rz(q1) =

1 0 0 0 c3 s3 0 −s3 c3

c2 0 −s2 0 1 0 s2 0 c2

c1 s1 0 −s1 c1 0 0 0 1

=

c1c2 s1c2 −s2 c1s2s3 −s1c3 s1s2s3 + c1c3 c2s3 c1s2c3 + s1s3 s1s2c3 −c1s3 c2c3

. (4.7) The formula for S can be obtained by expressing the joint velocity as a function of ˙q, and using S = ∂vJ/∂˙q. We shall skip the details, and just present the final result: S =

−s2 0 1 c2s3 c3 0 c2c3 −s3 0 0 0 0 0 0 0 0 0 0

. (4.8) The columns of S contain the coordinates in Fs of the axes about which the rotations take place: the first column is the z axis of Fp; the second is where the y axis was after the first rotation; and the third is where the x axis was after the first two rotations. As S is not a constant, we also need a formula for cJ. This is given by cJ = ˚ S ˙q =

−c2 ˙q2 0 0 −s2s3 ˙q2 + c2c3 ˙q3 −s3 ˙q3 0 −s2c3 ˙q2 −c2s3 ˙q3 −c3 ˙q3 0 0 0 0 0 0 0 0 0 0

˙q1 ˙q2 ˙q3

=

−c2 ˙q1 ˙q2 −s2s3 ˙q1 ˙q2 + c2c3 ˙q1 ˙q3 −s3 ˙q2 ˙q3 −s2c3 ˙q1 ˙q2 −c2s3 ˙q1 ˙q3 −c3 ˙q2 ˙q3 0 0 0

. (4.9) Equations 4.7 and 4.8 replace the formulae for E and S in Table 4.1, but the formula for T remains unchanged. This is therefore an example of a joint model in which S varies while T remains constant. Such a thing is possible because range(S) is a constant, and T is required only to satisfy range(T ) = range(S)⊥.

Euler Parameters A general rotation can also be described by means of a unit vector and an angle. The unit vector specifies the axis of rotation, and the angle specifies the magnitude. Euler parameters are closely related to this representation. If u and θ are the unit vector and angle, then the Euler parameters describing the same rotation are p0 = cos(θ/2) p1 = sin(θ/2)ux p2 = sin(θ/2)uy p3 = sin(θ/2)uz . (4.10) They are not independent, but must satisfy the equation p2 0 + p2 1 + p2 2 + p2 3 = 1 . (4.11) A spherical joint must have three independent velocity variables. Therefore, if we wish to use the four non-independent Euler parameters as position variables, then the velocity variables cannot be the derivatives of the position variables. The best choice for the velocity variables is to use the Cartesian coordinates in Fs of the joint’s angular velocity vector, ω. This choice results in S taking the simple, constant value shown in Table 4.1. We therefore need only a formula for E and a formula for the derivatives of the position variables. These are E = 2

p2 0 + p2 1 −1/2 p1p2 + p0p3 p1p3 −p0p2 p1p2 −p0p3 p2 0 + p2 2 −1/2 p2p3 + p0p1 p1p3 + p0p2 p2p3 −p0p1 p2 0 + p2 3 −1/2

(4.12) and

˙p0 ˙p1 ˙p2 ˙p3

= 1 2

−p1 −p2 −p3 p0 −p3 p2 p3 p0 −p1 −p2 p1 p0

ωx ωy ωz

. (4.13) A spherical joint model based on these equations is numerically superior to one based on Eqs. 4.7, 4.8 and 4.9. However, it does place a burden on the surrounding software of having to cope with a joint model in which the number of position variables differs from the number of velocity variables. Three more complications are: the values of the position variables must be obtained by integrating Eq. 4.13 rather than directly integrating the velocity variables; a zero rotation is defined by a nonzero set of Euler parameters; and it is necessary to renormalize Euler parameters frequently during the integration process. (Truncation errors during integration cause the parameters to fail to satisfy Eq. 4.11 exactly, and this causes a linear scaling error affecting the whole subtree supported by the joint. Ideally, they should be renormalized after every integration step.)

data for a extra data for a kinematic tree closed-loop system NB NL λ(1) · · · λ(NB) p(NB + 1) · · · p(NB + NL) jtype(1) · · · jtype(NB) s(NB + 1) · · · s(NB + NL) XT(1) · · · XT(NB) jtype(NB + 1) · · · jtype(NB + NL) I1 · · · INB XP(1) · · · XP(NL) XS(1) · · · XS(NL) Table 4.3: Complete system model 4.6 A Complete System Model Table 4.3 presents a complete system model for a rigid-body system. For a kinematic tree, the required data are: NB, the parent array, an array of joint type descriptors, the tree transform array, and an array of rigid-body inertias. The inertias are expressed in body coordinates, in which they are constants. For a closed-loop system, the required extra data are: NL and, for each loop joint, the predecessor and successor body numbers, the joint type descriptor, and the predecessor and successor transforms. In the simplest case, a software package might cater for a restricted set of joint types, such as revolute and prismatic joints only. In this case, jtype(i) could simply be a numeric type code, and it would be practical to hard-wire the appropriate values for XJ, vJ and S directly into the source code for jcalc, or even directly into the source code of a dynamics algorithm. Another simplification concerns access into vectors of joint variables, such as q and ˙q. In the general case, qi is a subvector of q, and its location within q depends on the sizes of all of the preceding subvectors. Thus, the code that accesses and updates qi should have an array of offsets, one per joint, indicating which element of q is the first element of qi. However, if every joint has exactly one degree of freedom, as is the case for revolute and prismatic joints, then qi is simply element i of q. If the objective is to develop a general-purpose software package, then it makes more sense to create a library of joint models, and have jcalc access this library. One particularly useful kind of joint is a user-defined joint. This facility could be provided by making the library extensible. Alternatively, one could simply require that one of the ‘parameters’ of a user-defined joint is the address of a user-supplied jcalc for that joint, and the built-in jcalc could be programmed to pass control to the user-supplied function.
